\documentclass[a4paper,fontset=none]{ctexbook}

\usepackage{anyfontsize}
\usepackage{amsmath,mathrsfs,wasysym}
\usepackage{graphicx,caption,subcaption}
\captionsetup{font=small}
\usepackage{geometry,parskip}
\geometry{top=3cm,bottom=3cm,left=2.75cm,right=2.75cm}
\usepackage{indentfirst}
\setlength{\parindent}{0em}
\usepackage[T1]{fontenc}
\usepackage[libertine,vvarbb]{newtxmath}
\usepackage[scr=rsfso,frak=euler,bb=ams]{mathalfa}
\usepackage{bm}
\usepackage{fontspec}
\setmainfont{Linux Libertine O}
\setsansfont{Linux Biolinum O}
\setCJKmainfont{SimSun}[BoldFont=Noto Sans CJK SC Medium]

\begin{document}

    \title{\textbf{实验三十八\ \ 微波的布拉格衍射}}
    \author{1900011604 张植竣}
    \date{2021年4月24日}

    \maketitle

    \section*{1\ \ 实验数据和现象}
    
    \subsection*{1.1\ \ 微波的布拉格衍射}
    
    \textbf{[100]晶面}
    
    实验数据如下：
    
    \begin{table}[h!]
        \begin{center}
            \begin{tabular}{|c|c|c|c|c|c|c|c|c|c|}
                \hline
                $\beta/^\circ$ & 20 & 25 & 30 & 32 & 33 & 34 & 35 & 36 & 37 \\
                $I/\mu\mathrm{A}$ & 0 & 1 & 4 & 3 & 6 & 10 & 18 & 24 & 24 \\
                \hline
                $\beta/^\circ$ & 38 & 39 & 40 & 42 & 45 & 50 & 55 & 60 & 62 \\
                $I/\mu\mathrm{A}$ & 24 & 26 & 21 & 16 & 0 & 0 & 0 & 3 & 6 \\
                \hline
                $\beta/^\circ$ & 64 & 65 & 66 & 66.5 & 67 & 67.5 & 68 & 68.5 & 69 \\
                $I/\mu\mathrm{A}$ & 18 & 12 & 18 & 20 & 45 & 52 & 63 & 57 & 52 \\
                \hline
                $\beta/^\circ$ & 69.5 & 70 & 70.5 & 71 & 71.5 & 72 & 72.5 & 73 & 73.5 \\
                $I/\mu\mathrm{A}$ & 33 & 29 & 24 & 30 & 46 & 64 & 67 & 61 & 38 \\
                \hline
                $\beta/^\circ$ & 74 & 75 & 78 & 80 & 81 & 82 & 83 & 84 & 85 \\
                $I/\mu\mathrm{A}$ & 14 & 4 & 30 & 22 & 32 & 20 & 14 & 34 & 70 \\
                \hline
            \end{tabular}
        \end{center}
        \caption{[100]晶面的实验数据}
    \end{table}
    
    绘制模拟晶体[100]晶面的布拉格衍射曲线如下：
    
    \begin{figure}[!ht]
        \centering
        \includegraphics[width=0.75\textwidth]{1.png}
        \caption{[100]晶面的布拉格衍射曲线}
    \end{figure}
    
    由上图修正后的蓝色曲线测量得到衍射极大处入射角为$\beta_1 = 70.26^\circ$，$\beta_2 = 38.44^\circ$。
    
    取$f = 9370\,\mathrm{MHz}$（对应$\lambda = 3.199492615\,\mathrm{cm}$），晶格常数$a = 4.0\,\mathrm{cm}$，由公式$2d\cos\beta = k\lambda$计算出衍射极大处入射角的理论值为：
    \[
    \beta_1^* = \arccos\left(\frac{\lambda}{2d}\right) = 66.4258^\circ,\qquad \beta_2^* = \arccos\left(\frac{\lambda}{2d}\right) = 36.8820^\circ
    \]
    实验值与理论值的差异为：
    \[
    \Delta\beta_1 = \beta_1 - \beta_1^* = 3.83^\circ,\qquad \Delta\beta_2 = \beta_2 - \beta_2^* = 1.56^\circ
    \]
    
    \bigskip
    \textbf{[110]晶面}
    
    实验数据如下：
    
    \begin{table}[h!]
        \begin{center}
            \begin{tabular}{|c|c|c|c|c|c|c|c|c|c|c|}
                \hline
                $\beta/^\circ$ & 20 & 25 & 30 & 35 & 40 & 45 & 47 & 49 & 49.5 & 50 \\
                $I/\mu\mathrm{A}$ & 4 & 4 & 0 & 1 & 3 & 6 & 3 & 7 & 9 & 15 \\
                \hline
                $\beta/^\circ$ & 50.5 & 51 & 51.5 & 52 & 52.5 & 53 & 53.5 & 54 & 54.5 & 55 \\
                $I/\mu\mathrm{A}$ & 21 & 32 & 20 & 36 & 32 & 34 & 28 & 26 & 18 & 10 \\
                \hline
                $\beta/^\circ$ & 55.5 & 56 & 57 & 58 & 59 & 60 & 65 & 70 & 75 & 78 \\
                $I/\mu\mathrm{A}$ & 3 & 1 & 1 & 1 & 2 & 4 & 16 & 34 & 33 & 57 \\
                \hline
            \end{tabular}
        \end{center}
        \caption{[110]晶面的实验数据}
    \end{table}
    
    绘制模拟晶体[110]晶面的布拉格衍射曲线如下：
    
    \begin{figure}[!ht]
        \centering
        \includegraphics[width=0.75\textwidth]{2.png}
        \caption{[110]晶面的布拉格衍射曲线}
    \end{figure}
    
    由上图修正后的蓝色曲线测量得到衍射极大处入射角为$\beta_1 = 51.97^\circ$。
    
    取$f = 9370\,\mathrm{MHz}$（对应$\lambda = 3.199492615\,\mathrm{cm}$），晶格常数$a = 4.0\,\mathrm{cm}$，由公式$2d\cos\beta = k\lambda$计算出衍射极大处入射角的理论值为：
    \[
    \beta_1^* = \arccos\left(\frac{\lambda}{2d}\right) = 55.5563^\circ
    \]
    实验值与理论值的差异为：
    \[
    \Delta\beta_1 = \beta_1 - \beta_1^* = -3.59^\circ
    \]

    \subsection*{1.2\ \ 微波的单缝衍射}
    
    实验数据如下：（其中$\theta = 0^\circ$时$I = 96\,\mu\mathrm{A}$）
    
    \begin{table}[h!]
        \begin{center}
            \begin{tabular}{|c|c|c|c|c|c|c|}
                \hline
                $\theta/^\circ$ & 5 & 10 & 15 & 20 & 25 & 26 \\
                $I/\mu\mathrm{A}$ & 92 & 80 & 44 & 13 & 1 & 0.5 \\
                \hline
                $\theta/^\circ$ & 27 & 28 & 29 & 30 & 35 & 40 \\
                $I/\mu\mathrm{A}$ & 0 & 0 & 0 & 0 & 1 & 1.5 \\
                \hline
                $\theta/^\circ$ & 45 & 50 & 55 & -55 & -50 & -45 \\
                $I/\mu\mathrm{A}$ & 2 & 1.5 & 1 & 0 & 1 & 0.5 \\
                \hline
                $\theta/^\circ$ & -40 & -35 & -30 & -29 & -28 & -27 \\
                $I/\mu\mathrm{A}$ & 0.5 & 1 & 0 & 0 & 0.5 & 0.5 \\
                \hline
                $\theta/^\circ$ & -26 & -25 & -20 & -15 & -10 & -5 \\
                $I/\mu\mathrm{A}$ & 1 & 1.5 & 18 & 48 & 73 & 94 \\
                \hline
            \end{tabular}
        \end{center}
        \caption{微波的单缝衍射}
    \end{table}
    
    \begin{figure}[!ht]
        \centering
        \includegraphics[width=0.75\textwidth]{3.png}
        \caption{微波的单缝衍射}
    \end{figure}
    
    根据拟合曲线测得左侧谷值对应角度$\theta_1 = -29.43^\circ$，右侧谷值对应角度$\theta_2 = 27.37^\circ$，测量得到$a = 7.0\,\mathrm{cm}$，计算波长：
    \[
    \lambda = a\sin{\overline\theta} = a\sin\left(\frac{\theta_2-\theta_1}{2}\right) = 3.329369463\,\mathrm{cm}
    \]
    根据贝塞尔公式估计角度的测量不确定度为$\sigma_\theta = 1.456639969^\circ$，又估计单缝宽度的不确定度为$\sigma_a = 0.1\,\mathrm{cm}$，则测量得到的波长$\lambda$的不确定度为：
    \[
    \sigma_\lambda = \sqrt{\left(\sigma_a\sin{\overline\theta}\right)^2 + 2\times\left(\frac{1}{2}a\sigma_\theta\cos{\overline\theta}\right)^2} = 0.120479124
    \]
    即测得波长为：
    \[
    \lambda = (3.33\pm0.12)\,\mathrm{cm}
    \]
    与理论值$\lambda^*=3.20\,\mathrm{cm}$较为符合。

    \subsection*{1.3\ \ 微波的迈克耳孙干涉}
    
    观察电表记录接收信号的变化并按顺序记录出现干涉极小值时B板的位置$x_{\mathrm{min}}$和干涉极大值时B板的位置$x_{\mathrm{max}}$，分别对$x_{\mathrm{min}}$和$x_{\mathrm{max}}$做最小二乘法拟合，并由$x_i - i$的斜率求出微波的波长$\lambda$。
    
    \textbf{干涉极小值}
    
    出现干涉极小值时B板的位置：
    
    \begin{table}[h!]
        \begin{center}
            \begin{tabular}{|c|c|c|c|c|c|}
                \hline
                $i$ & 1 & 2 & 3 & 4 \\
                \hline
                $x_{\mathrm{min}}/\mathrm{mm}$ & 11.0 & 27.5 & 45.3 & 59.8 \\
                \hline
            \end{tabular}
        \end{center}
        \caption{出现干涉极小值时B板的位置}
    \end{table}
    
    最小二乘法拟合得到：
    
    \begin{table}[h!]
        \begin{center}
            \begin{tabular}{|c|c|c|c|c|c|}
                \hline
                $k$ & $r$ & $\sigma_{k,A}$ & $\sigma_{k,B}$ & $\sigma_k$ \\
                \hline
                16.42 & 0.999237566 & 0.453651849 & 0.025819889 & 0.454386033 \\
                \hline
            \end{tabular}
        \end{center}
    \end{table}
    
    其中：（$n = 4$，估计$e_x = 0.1\,\mathrm{mm}$，下同）
    \[
    \sigma_{k,A} = k\sqrt{\frac{1/r^2 - 1}{n - 2}},\qquad
    \sigma_{k,B} = \frac{e_x/\sqrt{3}}{\sqrt{\sum\limits_i^n(i - \overline{i})^2}},\qquad
    \sigma_k = \sqrt{\sigma_{k,A}^2 + \sigma_{k,B}^2}
    \]
    又因为
    \[
    \lambda = \dfrac{k}{5},\qquad \sigma_\lambda = \dfrac{\sigma_k}{5}
    \]
    测得微波的波长为：
    \[
    \lambda = (3.28\pm0.09)\,\mathrm{cm}
    \]
    
    \bigskip
    \textbf{干涉极大值}
    
    出现干涉极大值时B板的位置：
    
    \begin{table}[h!]
        \begin{center}
            \begin{tabular}{|c|c|c|c|c|c|}
                \hline
                $i$ & 1 & 2 & 3 & 4 \\
                \hline
                $x_{\mathrm{max}}/\mathrm{mm}$ & 11.0 & 27.5 & 45.3 & 59.8 \\
                \hline
            \end{tabular}
        \end{center}
        \caption{出现干涉极大值时B板的位置}
    \end{table}
    
    最小二乘法拟合得到：
    
    \begin{table}[h!]
        \begin{center}
            \begin{tabular}{|c|c|c|c|c|c|}
                \hline
                $k$ & $r$ & $\sigma_{k,A}$ & $\sigma_{k,B}$ & $\sigma_k$ \\
                \hline
                15.96 & 0.999936407 & 0.127279220 & 0.025819889 & 0.129871732 \\
                \hline
            \end{tabular}
        \end{center}
    \end{table}
    
    测得微波的波长为：
    \[
    \lambda' = (3.192\pm0.026)\,\mathrm{cm}
    \]
    
    \bigskip
    \textbf{与理论值对比}
    
    理论值为$\lambda^* = \dfrac{c}{f} = 3.199492615\,\mathrm{cm}$，与$\lambda'$更接近，可见用干涉极大值测量微波波长误差更小。

    \section*{2\ \ 分析与讨论}
    
    1. 本实验中使用螺旋测微器调节频率，导致波长的理论值存在误差，但是相对误差在$10^{-3}$量级，远小于测量误差，仍然可以认为$\lambda^* = 3.20\,\mathrm{cm}$。
        
    2. 在测量微波的布拉格衍射时，会发现衍射峰分裂为多个峰，这里采用的是去除异常下降的值之后拟合光滑曲线寻找原本衍射峰的位置，会带来一定误差。
    
    3. 微波波长较长，传播方向、强度受外界干扰较大，使得测量精度相比可见光明显降低，误差可以达到几度。
\end{document}
